\chapter{John Willard Milnor --- Fields Medal 1962}\label{ch1962johnwillardmilnor:john-willard-milnor-fields-medal-1962-a-comprehensive-report}

\begin{center}\rule{0.5\linewidth}{0.5pt}\end{center}

\textbf{Prepared:} July 2026 \textbf{Subject:} John Willard Milnor (b.
February 20, 1931, Orange, New Jersey, USA) \textbf{Position:}
Professor, Princeton University; SUNY Stony Brook \textbf{Doctoral
Advisor:} Ralph H. Fox (Princeton University) \textbf{Fields Medal
Citation:} \emph{``Proved that a 7-dimensional sphere can have several
differential structures; this led to the creation of the field of
differential topology.''}

\begin{center}\rule{0.5\linewidth}{0.5pt}\end{center}

\section{1. Introduction}\label{ch1962johnwillardmilnor:introduction}

John Willard Milnor was awarded the Fields Medal in 1962 at age 31 for
his substantial work in differential topology, well-knownly his
discovery of exotic spheres --- differentiable manifolds homeomorphic
but not diffeomorphic to the standard 7-sphere. This result inaugurated
differential topology and changed how we think about
smooth structures.

Milnor's mathematical range is notable. He has made fundamental
contributions to topology, differential geometry, algebraic K-theory,
dynamics, and knot theory. His monographs --- \emph{Topology from the
Differentiable Viewpoint}, \emph{Morse Theory}, \emph{Characteristic
Classes}, and \emph{Lectures on the h-Cobordism Theorem} --- are
classics that have trained generations of mathematicians.

Major contributions include: - \textbf{Exotic spheres}: \(S^7\) admits
28 distinct differentiable structures, founding differential topology. -
\textbf{Milnor K-theory}: Algebraic K-theory for fields connecting to
Galois cohomology and motivic cohomology. - \textbf{The Milnor fibration
and Milnor number}: Invariants for isolated hypersurface singularities.
- \textbf{Milnor's theorem on total curvature of knots}: Relating
curvature to knot type. - \textbf{Characteristic classes and classifying
spaces}. - \textbf{Holomorphic dynamics}: Classification of rational
maps. - \textbf{Milnor's theorem on growth of groups}.

\begin{center}\rule{0.5\linewidth}{0.5pt}\end{center}

\section{2. Biographical Background}\label{ch1962johnwillardmilnor:biographical-background}

\subsection{2.1 Early Life and
Education}\label{ch1962johnwillardmilnor:early-life-and-education}

John Willard Milnor was born on February 20, 1931, in Orange, New
Jersey. His father was an engineer and his mother a schoolteacher. He
showed exceptional mathematical talent early on, winning a scholarship
to Princeton University at age 16. As an undergraduate at Princeton, he
became a Putnam Fellow in 1949, scoring among the top five in the
William Lowell Putnam Competition. He continued at Princeton for
graduate studies under Ralph Fox, completing his Ph.D.~in 1954.

\subsection{2.2 Academic Career}\label{ch1962johnwillardmilnor:academic-career}

\begin{itemize}
\tightlist
\item
 \textbf{1954:} Ph.D., Princeton University, under Ralph H. Fox.
 Thesis: ``Isotopy of Links.''
\item
 \textbf{1954--1960:} Faculty, Princeton University.
\item
 \textbf{1960--1962:} Professor, Princeton University (Henry Putman
 chair from 1962).
\item
 \textbf{1962:} Fields Medal at the ICM in Stockholm.
\item
 \textbf{1968--1970:} Professor, Massachusetts Institute of Technology.
\item
 \textbf{1970--1988:} Professor, Institute for Advanced Study,
 Princeton.
\item
 \textbf{1988--2003:} Professor, SUNY Stony Brook.
\item
 \textbf{1989:} Wolf Prize in Mathematics.
\item
 \textbf{2011:} Abel Prize.
\item
 \textbf{Present:} Professor Emeritus, SUNY Stony Brook.
\end{itemize}

\subsection{2.3 Doctoral Work}\label{ch1962johnwillardmilnor:doctoral-work}

Milnor's Ph.D.~thesis on the isotopy of links already showed his
characteristic originality. He studied when two links in \(S^3\) are
isotopic (i.e., can be deformed into each other without cutting) and
introduced new invariants for distinguishing knots and links. This work
foreshadowed his lifelong interest in the classification of geometric
objects.

\begin{center}\rule{0.5\linewidth}{0.5pt}\end{center}

\section{3. Exotic Spheres and Differential
Topology}\label{ch1962johnwillardmilnor:exotic-spheres-and-differential-topology}

\subsection{3.1 Background}\label{ch1962johnwillardmilnor:background}

Until Milnor's work, it was widely believed that a given topological
manifold could possess at most one differentiable structure. The
7-sphere \(S^7\) seemed particularly well understood: it is the set of
unit vectors in \(\mathbb{R}^8\), with the standard smooth structure
inherited from \(\mathbb{R}^8\).

\subsection{3.2 The Discovery}\label{ch1962johnwillardmilnor:the-discovery}

In a landmark 1956 paper ``On Manifolds Homeomorphic to the 7-Sphere,''
Milnor constructed a 7-manifold \(M^7\) that is homeomorphic to \(S^7\)
but not diffeomorphic to it. He considered 3-sphere bundles over \(S^4\)
with structure group \(SO(4)\). Identifying quaternions \(\mathbb{H}\)
with \(\mathbb{R}^4\), a map \(f: S^3 \to SO(4)\) is given by
\(f(u)v = u^h v u^j\) where \(h+j=1\) and \(h-j=k\) is an odd integer.
The total space \(M_k\) of the associated bundle is homeomorphic to
\(S^7\) for all odd \(k\).

Milnor defined an invariant \(\lambda(M^7) \in \mathbb{Z}/7\mathbb{Z}\)
using Thom's cobordism theory. If \(M^7 = \partial B^8\), then
\(\lambda(M) = 2q(B) - \tau(B) \mod 7\), where \(q(B)\) is a Pontryagin
number and \(\tau(B)\) is the signature. He computed
\(\lambda(M_k) \equiv k^2-1 \mod 7\), so for
\(k \not\equiv \pm 1 \mod 7\), \(\lambda \neq 0\), proving these are
exotic spheres.

\subsection{3.3 Classification}\label{ch1962johnwillardmilnor:classification}

With Michel Kervaire, Milnor classified exotic spheres in all
dimensions. The group \(\Theta_n\) of homotopy spheres (manifolds
homotopy equivalent to \(S^n\)) under connected sum is: - Finite cyclic
for \(n \neq 4\). - For \(n=7\), \(\Theta_7 \cong \mathbb{Z}/28\), so
there are 28 distinct differentiable structures on \(S^7\).

The classification uses the Adams spectral sequence, the J-homomorphism
of George Whitehead, and deep results from algebraic topology. The order
of \(\Theta_n\) is related to Bernoulli numbers via the Todd
polynomials, reflecting a deep connection between differential topology
and number theory.

\subsection{3.4 Impact}\label{ch1962johnwillardmilnor:impact}

Exotic spheres founded differential topology. They show that the
category of topological manifolds differs substantially from the
category of smooth manifolds. The work led to the development of surgery
theory by Browder, Novikov, Wall, and others, and to the complete
classification of high-dimensional manifolds. The discovery also has
implications for mathematical physics, where exotic spheres appear in
string theory and gauge theory.

\begin{center}\rule{0.5\linewidth}{0.5pt}\end{center}

\section{4. The h-Cobordism Theorem}\label{ch1962johnwillardmilnor:the-h-cobordism-theorem}

\subsection{4.1 Background}\label{ch1962johnwillardmilnor:background-1}

A cobordism between two manifolds \(M_0\) and \(M_1\) is a manifold
\(W\) with boundary \(\partial W = M_0 \cup M_1\). It is an
\textbf{h-cobordism} if the inclusions \(M_0 \hookrightarrow W\) and
\(M_1 \hookrightarrow W\) are homotopy equivalences.

\subsection{4.2 The Smale--Milnor
Theorem}\label{ch1962johnwillardmilnor:the-smalemilnor-theorem}

Stephen Smale proved the h-cobordism theorem for \(n \geq 5\): if \(W\)
is a simply connected h-cobordism between \(n\)-dimensional manifolds,
then \(W \cong M_0 \times [0,1]\), so \(M_0 \cong M_1\). Milnor wrote
the definitive account of this theorem in \emph{Lectures on the
h-Cobordism Theorem}, providing the complete Morse-theoretic proof.

The proof uses \textbf{Morse theory} on the cobordism to find a function
\(f: W \to [0,1]\) with no critical points. Milnor's exposition
clarified the handlebody decomposition, the Whitney trick for removing
intersections, and the fundamental group considerations.

\subsection{4.3 The Generalized Poincaré
Conjecture}\label{ch1962johnwillardmilnor:the-generalized-poincaruxe9-conjecture}

The h-cobordism theorem implies the \textbf{Generalized Poincaré
Conjecture} for \(n \geq 5\): any homotopy \(n\)-sphere is homeomorphic
to \(S^n\). Combined with the Kervaire--Milnor classification of exotic
spheres, this shows that for \(n \geq 5\), the set of differentiable
structures on \(S^n\) is finite but can be non-trivial.

\begin{center}\rule{0.5\linewidth}{0.5pt}\end{center}

\section{5. Milnor K-Theory}\label{ch1962johnwillardmilnor:milnor-k-theory}

\subsection{5.1 Definition}\label{ch1962johnwillardmilnor:definition}

For a field \(F\), the \textbf{Milnor K-theory} \(K_*^M(F)\) is the
graded commutative algebra generated by symbols \(\{a_1,\ldots,a_n\}\)
for \(a_i \in F^\times\), subject to: - \textbf{Bilinearity}:
\(\{a,b\} = \{ab\} - \{a\} - \{b\}\). - \textbf{Steinberg relation}:
\(\{a,1-a\} = 0\) for \(a \neq 0,1\).

Thus \(K_1^M(F) \cong F^\times\), and \(K_2^M(F)\) is generated by
symbols \(\{a,b\}\) with the Steinberg relation.

\subsection{5.2 Connections to Galois
Cohomology}\label{ch1962johnwillardmilnor:connections-to-galois-cohomology}

Milnor made the notable conjecture that for any field \(F\) and
integer \(n\), \(K_n^M(F)/2 K_n^M(F)\) is isomorphic to the Galois
cohomology group \(H^n(F, \mathbb{Z}/2\mathbb{Z})\). This was proved by
Voevodsky (for which he won the Fields Medal), using motivic cohomology
and the norm residue homomorphism. The generalization to \(\ell\)-adic
coefficients (the Bloch--Kato conjecture) was also proved by Voevodsky.

\subsection{5.3 Applications}\label{ch1962johnwillardmilnor:applications}

Milnor K-theory is a fundamental tool in: - \textbf{Algebraic K-theory}:
\(K_2\) of fields is described by Matsumoto's theorem in terms of Milnor
K-theory. - \textbf{Quadratic forms}: The Milnor conjectures relate to
the Witt ring of quadratic forms. - \textbf{Class field theory}: Higher
class field theory uses Milnor K-theory to describe abelian extensions.
- \textbf{Motivic cohomology}: Milnor K-theory is the cohomology of the
motivic complex \(\mathbb{Z}(n)\) at the \(n\)th slice.

\begin{center}\rule{0.5\linewidth}{0.5pt}\end{center}

\section{6. The Milnor Fibration and the Milnor
Number}\label{ch1962johnwillardmilnor:the-milnor-fibration-and-the-milnor-number}

\subsection{6.1 The Milnor Fibration}\label{ch1962johnwillardmilnor:the-milnor-fibration}

For a polynomial \(f: \mathbb{C}^{n+1} \to \mathbb{C}\) with an isolated
singularity at \(0\), consider the sphere \(S_\epsilon\) of radius
\(\epsilon\) around \(0\). The map \(\phi(z) = f(z)/|f(z)|\) gives a
fibration:

\[f: S_\epsilon \setminus f^{-1}(0) \to S^1\]

known as the \textbf{Milnor fibration}. The fiber
\(F = f^{-1}(1) \cap B_\epsilon\) is a smooth manifold of dimension
\(2n\), called the \textbf{Milnor fiber}. It has the homotopy type of a
bouquet of \(n\)-spheres.

\subsection{6.2 The Milnor Number}\label{ch1962johnwillardmilnor:the-milnor-number}

The \textbf{Milnor number}
\(\mu(f) = \dim_\mathbb{C} \mathcal{O}_n / J_f\), where
\(J_f = \langle \partial f/\partial z_1, \ldots, \partial f/\partial z_{n+1} \rangle\)
is the Jacobian ideal. It equals the number of \(n\)-spheres in the
bouquet forming the Milnor fiber: \(F \simeq \bigvee_{\mu} S^n\).

Milnor's formula relates the Milnor number to other invariants:

\[\mu = \dim H_n(F, \mathbb{C}) = \tau(f) + \sigma(f)\]

where \(\tau(f)\) is the Tjurina number and \(\sigma(f)\) is related to
the signature of the intersection form.

\subsection{6.3 Impact}\label{ch1962johnwillardmilnor:impact-1}

The Milnor fibration and Milnor number are central to singularity
theory. They connect algebraic geometry (through the Jacobian ring),
topology (through the vanishing cycles), and differential geometry
(through the intersection form). Sebastiani--Thom type theorems express
the Milnor number of a sum of polynomials in terms of the summands, and
the theory underlies the study of deformations of hypersurface
singularities.

\begin{center}\rule{0.5\linewidth}{0.5pt}\end{center}

\section{7. Curvature of Knots}\label{ch1962johnwillardmilnor:curvature-of-knots}

\subsection{7.1 Milnor's Theorem}\label{ch1962johnwillardmilnor:milnors-theorem}

For a smooth knot \(K \subset \mathbb{R}^3\), the \textbf{total
curvature} is \(\int_K |\kappa| ds\), where \(\kappa\) is the curvature.
Milnor proved a beautiful theorem:

\[\int_K |\kappa| \, ds \geq 4\pi\]

Equality holds if and only if \(K\) is a convex planar curve (a trivial
knot lying in a plane).

This result connects the geometry of a knot (its curvature) to its
topology (its knot type). The proof uses the \textbf{Fáry--Milnor
theorem}: a knot in \(\mathbb{R}^3\) with total curvature less than
\(4\pi\) must be unknotted.

\subsection{7.2 The Fáry--Milnor
Theorem}\label{ch1962johnwillardmilnor:the-fuxe1rymilnor-theorem}

The Fáry--Milnor theorem states that if the total curvature of a closed
curve in \(\mathbb{R}^3\) is less than \(4\pi\), then the curve is
unknotted. It uses the Gauss map \(T: S^1 \to S^2\) (the unit tangent
vector), whose image area is bounded by the total curvature. If the area
is less than \(2\pi\) (i.e., total curvature \(< 4\pi\)), the Gauss map
is not surjective, and there exists a direction from which the knot
projects to a planar curve with no crossings.

\begin{center}\rule{0.5\linewidth}{0.5pt}\end{center}

\section{8. Morse Theory}\label{ch1962johnwillardmilnor:morse-theory}

\subsection{8.1 Background}\label{ch1962johnwillardmilnor:background-2}

\textbf{Morse theory} studies the topology of a manifold through a
smooth function \(f: M \to \mathbb{R}\) with non-degenerate critical
points. The \textbf{Morse inequalities} relate the numbers \(c_k\) of
critical points of index \(k\) to the Betti numbers \(b_k\):
\(c_k - c_{k-1} + \cdots \pm c_0 \geq b_k - b_{k-1} + \cdots \pm b_0\).

\subsection{8.2 Milnor's Contributions}\label{ch1962johnwillardmilnor:milnors-contributions}

Milnor wrote the classic text \emph{Morse Theory} based on his 1963
lectures at Princeton, providing the definitive treatment of the
fundamental theorems, the Lefschetz hyperplane theorem, Bott periodicity
via Morse theory on the loop space, and applications to geodesics. His
proof of Bott periodicity using Morse theory remains one of the most
elegant in topology.

\begin{center}\rule{0.5\linewidth}{0.5pt}\end{center}

\section{9. Characteristic Classes}\label{ch1962johnwillardmilnor:characteristic-classes}

\subsection{9.1 Stiefel--Whitney and Pontryagin
Classes}\label{ch1962johnwillardmilnor:stiefelwhitney-and-pontryagin-classes}

\textbf{Characteristic classes} are cohomological invariants of vector
bundles that measure twisting. Milnor's book \emph{Characteristic
Classes} (with James Stasheff) is the standard reference, covering
Stiefel--Whitney classes \(w_i(E) \in H^i(B; \mathbb{Z}/2)\), Pontryagin
classes \(p_i(E) \in H^{4i}(B; \mathbb{Z})\), Chern classes
\(c_i(E) \in H^{2i}(B; \mathbb{Z})\), and the Euler class. The book
develops these through the axiomatic approach: naturality, the Whitney
sum formula, and normalization, allowing computation via the splitting
principle.

\subsection{9.3 The Milnor--Husemeller
Work}\label{ch1962johnwillardmilnor:the-milnorhusemeller-work}

With Dale Husemoller, Milnor studied symmetric bilinear forms over the
integers and their relation to characteristic classes and the signature
theorem, connecting quadratic forms to the topology of manifolds.

\begin{center}\rule{0.5\linewidth}{0.5pt}\end{center}

\section{10. The Classifying Space of a Topological
Group}\label{ch1962johnwillardmilnor:the-classifying-space-of-a-topological-group}

\subsection{10.1 Construction}\label{ch1962johnwillardmilnor:construction}

Milnor constructed the \textbf{classifying space} \(BG\) of a
topological group \(G\) as the geometric realization of the simplicial
space \(EG = G * G * \cdots\) (the infinite join). The resulting space
\(EG\) is contractible, and \(G\) acts freely on it, so \(BG = EG/G\).
The universal bundle \(EG \to BG\) has the property that any principal
\(G\)-bundle \(P \to X\) is classified by a map \(f: X \to BG\), unique
up to homotopy, with \(P \cong f^* EG\). Milnor's construction provides
a concrete model for \(BG\) for any topological group, not just discrete
groups, and leads to the Milnor--Moore theorem on Hopf algebras and the
structure of \(H_*(BG)\).

\begin{center}\rule{0.5\linewidth}{0.5pt}\end{center}

\section{11. The Milnor--Moore Theorem on Hopf
Algebras}\label{ch1962johnwillardmilnor:the-milnormoore-theorem-on-hopf-algebras}

\subsection{11.1 Statement}\label{ch1962johnwillardmilnor:statement}

With John C. Moore, Milnor proved a foundational theorem about the
structure of connected Hopf algebras: if \(A\) is a connected graded
commutative Hopf algebra over a field of characteristic 0, then \(A\) is
isomorphic (as an algebra) to the symmetric algebra on its space of
primitives:

\[A \cong S(\mathcal{P}(A))\]

where
\(\mathcal{P}(A) = \{a \in A: \Delta(a) = a \otimes 1 + 1 \otimes a\}\).
This theorem is fundamental in algebraic topology: the cohomology
\(H^*(G; \mathbb{Q})\) of a compact Lie group is an exterior algebra on
odd-degree generators, and the dual Steenrod algebra \(A_*\) has a
structure described explicitly by Milnor's generators \(\xi_i, \tau_i\):

\[\Delta(\xi_k) = \sum_{i=0}^k \xi_{k-i}^{p^i} \otimes \xi_i\]

\begin{center}\rule{0.5\linewidth}{0.5pt}\end{center}

\section{12. The Milnor Invariant of
Links}\label{ch1962johnwillardmilnor:the-milnor-invariant-of-links}

\subsection{12.1 Definition}\label{ch1962johnwillardmilnor:definition-1}

The \textbf{Milnor \(\mu\)-invariants} are higher-order linking numbers
for links in \(S^3\), generalizing the linking number
\(\mu(ij) = \text{lk}(K_i, K_j)\). Milnor invariants are defined using
the lower central series of the link group \(\pi_1(S^3 \setminus L)\).
They detect when a link is Milnor-trivial (bounds disjoint Seifert
surfaces). The invariant \(\mu(123)\) detects Borromean linking: the
Borromean rings have \(\mu(123) = \pm 1\) while all pairwise linking
numbers vanish. Milnor invariants are concordance invariants and appear
as the simplest examples of finite-type (Vassiliev) invariants.

\begin{center}\rule{0.5\linewidth}{0.5pt}\end{center}

\section{13. Growth of Groups}\label{ch1962johnwillardmilnor:growth-of-groups}

\subsection{13.1 Milnor's Problem}\label{ch1962johnwillardmilnor:milnors-problem}

For a finitely generated group \(G\), the \textbf{growth function}
\(\gamma(n)\) counts the number of group elements expressible as words
of length \(\leq n\). Milnor proved that if a finitely generated group
\(G\) has polynomial growth, then \(G\) is \textbf{virtually nilpotent}
(has a nilpotent subgroup of finite index). This was later generalized
by Gromov. Milnor also proved that the growth rate of the fundamental
group \(\pi_1(M)\) of a compact Riemannian manifold \(M\) with negative
sectional curvature is always exponential, connecting group theory,
geometry, and dynamics.

\begin{center}\rule{0.5\linewidth}{0.5pt}\end{center}

\section{14. Dynamics and Holomorphic
Dynamics}\label{ch1962johnwillardmilnor:dynamics-and-holomorphic-dynamics}

\subsection{14.1 Dynamics of Rational
Maps}\label{ch1962johnwillardmilnor:dynamics-of-rational-maps}

In later work, Milnor turned to holomorphic dynamics. His book
\emph{Dynamics in One Complex Variable} is the standard graduate text.
He studied: - \textbf{The Mandelbrot set}: its local connectivity,
hyperbolic components, and bifurcation structure. - \textbf{The
classification of rational maps}: postcritically finite maps, Thurston's
theorem. - \textbf{The Fatou and Julia sets}: properties of periodic
cycles, the classification of Fatou components.

\subsection{14.2 Milnor's Results in
Dynamics}\label{ch1962johnwillardmilnor:milnors-results-in-dynamics}

Milnor made several significant discoveries: - \textbf{The
Milnor--Thurston theorem} on kneading theory for unimodal maps of the
interval. - \textbf{The topology of the Mandelbrot set}: proof that
hyperbolic components are simply connected. - \textbf{The classification
of rational maps with a periodic critical point}: the relationship
between combinatorial types and parameter spaces.

\subsection{14.3 Impact}\label{ch1962johnwillardmilnor:impact-2}

Milnor's work in dynamics brought the same geometric clarity and depth
that characterized his earlier work in topology. He introduced new
combinatorial methods for studying the iteration of rational maps, and
his expository writing made the field accessible to a broad mathematical
audience.

\begin{center}\rule{0.5\linewidth}{0.5pt}\end{center}

\section{15. The Milnor Conjectures}\label{ch1962johnwillardmilnor:the-milnor-conjectures}

\subsection{15.1 The First Milnor
Conjecture}\label{ch1962johnwillardmilnor:the-first-milnor-conjecture}

Milnor conjectured that for a field \(F\),
\(K_n^M(F)/2K_n^M(F) \cong H^n(G_F, \mathbb{Z}/2)\), where \(G_F\) is
the absolute Galois group. This was proved by Voevodsky using motivic
cohomology and the norm residue homomorphism.

\subsection{15.2 The Second Milnor
Conjecture}\label{ch1962johnwillardmilnor:the-second-milnor-conjecture}

The second Milnor conjecture relates Milnor K-theory modulo 2 to the
Witt ring of quadratic forms. It states that the natural map
\(K_n^M(F)/2 \to I^n/I^{n+1}\) is an isomorphism, where \(I\) is the
fundamental ideal of the Witt ring. This was proved by Orlov, Vishik,
and Voevodsky. The conjecture has deep implications for the algebraic
theory of quadratic forms and the Kaplansky conjecture on the level of
fields.

\subsection{15.3 Significance}\label{ch1962johnwillardmilnor:significance}

The proof of the Milnor conjectures is one of the major achievements of
21st-century algebraic geometry. They unify Galois cohomology, algebraic
K-theory, and the theory of quadratic forms, and they motivated much of
the development of motivic cohomology by Voevodsky.

\begin{center}\rule{0.5\linewidth}{0.5pt}\end{center}

\section{16. Further Topics}\label{ch1962johnwillardmilnor:further-topics}

\subsection{16.1 The Milnor--Swan Theorem}\label{ch1962johnwillardmilnor:the-milnorswan-theorem}

With Richard Swan, Milnor studied the relationship between projective
modules over group rings and vector bundles over classifying spaces. The
\textbf{Milnor--Swan theorem} states that for a finite group \(G\), the
category of finitely generated projective modules over \(\mathbb{R}[G]\)
is equivalent to the category of real vector bundles over \(BG\).

\subsection{16.2 The Kervaire--Milnor
Formula}\label{ch1962johnwillardmilnor:the-kervairemilnor-formula}

The group \(\Theta_n\) of homotopy spheres fits into an exact sequence:

\[0 \to \Theta_n \to \pi_n(\text{PL/O}) \to \mathbb{Z}/2\]

The Kervaire--Milnor formula gives:

\[|\Theta_n| = \frac{2^{2n-2}(2^{n-1}-1)B_n}{n} \cdot \alpha_n\]

where \(B_n\) is the Bernoulli number and \(\alpha_n\) is a factor
depending on parity.

\subsection{16.3 The Pontryagin--Thom
Construction}\label{ch1962johnwillardmilnor:the-pontryaginthom-construction}

The \textbf{Pontryagin--Thom construction} establishes an isomorphism
between the cobordism ring \(\Omega_*\) of manifolds and the homotopy
groups of Thom spectra. Milnor used it to relate the signature of a
manifold to the \(L\)-genus and to compute the invariant \(\lambda(M)\)
for exotic spheres.

\begin{center}\rule{0.5\linewidth}{0.5pt}\end{center}

\section{Milnor's Books and Expository
Writing}\label{ch1962johnwillardmilnor:milnors-books-and-expository-writing}

Milnor is renowned for the clarity and elegance of his expository
writing. His books include: - \textbf{Topology from the Differentiable
Viewpoint} (1965): A concise introduction to differential topology. -
\textbf{Morse Theory} (1963): The definitive treatment based on his
Princeton lectures. - \textbf{Lectures on the h-Cobordism Theorem}
(1965): The complete proof of the h-cobordism theorem. -
\textbf{Characteristic Classes} (with J. Stasheff, 1974): The standard
reference on characteristic classes. - \textbf{Symmetric Bilinear Forms}
(with D. Husemoller, 1973): Algebraic theory of quadratic forms. -
\textbf{Dynamics in One Complex Variable} (1999): The standard text on
holomorphic dynamics.

\begin{center}\rule{0.5\linewidth}{0.5pt}\end{center}

\section{Awards and Recognition}\label{ch1962johnwillardmilnor:awards-and-recognition}

\textbf{1962:} - Fields Medal (International Mathematical Union).

\textbf{1967:} - National Medal of Science (USA).

\textbf{1970:} - Lester R. Ford Award for mathematical exposition.

\textbf{1982:} - Leroy P. Steele Prize for Seminal Contribution to
Research.

\textbf{1989:} - Wolf Prize in Mathematics.

\textbf{2004:} - Leroy P. Steele Prize for Mathematical Exposition.

\textbf{2011:} - Abel Prize (Norwegian Academy of Science and Letters).

\textbf{Elected to:} - National Academy of Sciences (USA) (1963). -
American Academy of Arts and Sciences. - Royal Society of London,
Foreign Member. - Russian Academy of Sciences, Foreign Member. -
Norwegian Academy of Science and Letters. - London Mathematical Society,
Honorary Member.

\textbf{2022:} - Honorary doctorate, University of Chicago.

\begin{center}\rule{0.5\linewidth}{0.5pt}\end{center}

\section{Selected Major
Publications}\label{ch1962johnwillardmilnor:selected-major-publications}

\begin{enumerate}
\def\labelenumi{\arabic{enumi}.}
\item
 \textbf{``On manifolds homeomorphic to the 7-sphere.''} \emph{Annals
 of Mathematics} 64, 399--405, 1956. --- \emph{The discovery of exotic
 spheres, founding differential topology.}
\item
 \textbf{``Groups of homotopy spheres I.''} (with M. Kervaire)
 \emph{Annals of Mathematics} 77, 504--537, 1963. ---
 \emph{Classification of exotic spheres in all dimensions.}
\item
 \textbf{``Algebraic K-theory and quadratic forms.''} \emph{Inventiones
 Mathematicae} 9, 318--344, 1970. --- \emph{Introduction of Milnor
 K-theory.}
\item
 \textbf{``Singular points of complex hypersurfaces.''} \emph{Annals of
 Mathematics Studies} 61, Princeton, 1968. --- \emph{The Milnor
 fibration and Milnor number.}
\item
 \textbf{``On the total curvature of knots.''} \emph{Annals of
 Mathematics} 52, 248--257, 1950. --- \emph{Milnor's theorem on total
 curvature.}
\item
 \textbf{``Construction of universal bundles I--II.''} \emph{Annals of
 Mathematics} 63, 272--284; 63, 430--436, 1956. --- \emph{Classifying
 spaces of topological groups.}
\item
 \textbf{``On the structure of Hopf algebras.''} (with J. C. Moore)
 \emph{Annals of Mathematics} 81, 211--264, 1965. --- \emph{The
 Milnor--Moore theorem.}
\item
 \textbf{``Morse theory.''} \emph{Annals of Mathematics Studies} 51,
 Princeton, 1963. --- \emph{The definitive text on Morse theory.}
\item
 \textbf{``Characteristic classes.''} (with J. Stasheff) \emph{Annals
 of Mathematics Studies} 76, Princeton, 1974. --- \emph{The standard
 reference on characteristic classes.}
\item
 \textbf{``Dynamics in one complex variable.''} \emph{Vieweg}, 1999;
 3rd ed.~Princeton, 2006. --- \emph{The standard text on holomorphic
 dynamics.}
\item
 \textbf{``Topology from the differentiable viewpoint.''}
 \emph{University of Virginia Press}, 1965; Princeton, 1997. ---
 \emph{Classic introduction to differential topology.}
\item
 \textbf{``On the relationship between the growth of groups and
 curvature.''} \emph{Journal of Differential Geometry} 2, 1--6, 1968.
 --- \emph{Growth of groups and virtual nilpotency.}
\end{enumerate}

\begin{center}\rule{0.5\linewidth}{0.5pt}\end{center}

\section{References and Further
Reading}\label{ch1962johnwillardmilnor:references-and-further-reading}

\subsection{Official Sources}\label{ch1962johnwillardmilnor:official-sources}

\begin{itemize}
\tightlist
\item
 IMU Fields Medal 1962:
 https://www.mathunion.org/imu-awards/fields-medal/fields-medals-1962
\item
 Abel Prize 2011: https://www.abelprize.no/abelprize/2011/
\item
 Wolf Prize 1989: https://wolffund.org.il/john-milnor/
\end{itemize}

\subsection{Key Technical References}\label{ch1962johnwillardmilnor:key-technical-references}

\begin{itemize}
\tightlist
\item
 Milnor, J. \emph{Collected Papers} (3 vols.). Publish or Perish, 1994.
\item
 Milnor, J. \emph{Topology from the Differentiable Viewpoint}.
 Princeton University Press, 1997.
\item
 Milnor, J. \emph{Morse Theory}. Princeton University Press, 1963.
\item
 Milnor, J. and Stasheff, J. \emph{Characteristic Classes}. Princeton
 University Press, 1974.
\item
 Milnor, J. \emph{Dynamics in One Complex Variable}. Princeton
 University Press, 2006.
\end{itemize}

\subsection{Textbooks and Surveys}\label{ch1962johnwillardmilnor:textbooks-and-surveys}

\begin{itemize}
\tightlist
\item
 Hirsch, M. \emph{Differential Topology}. Springer, 1976.
\item
 Kosinski, A. \emph{Differential Manifolds}. Academic Press, 1993.
\item
 Ranicki, A. \emph{Algebraic and Geometric Surgery}. Oxford University
 Press, 2002.
\end{itemize}

\subsection{Biography and Exposition}\label{ch1962johnwillardmilnor:biography-and-exposition}

\begin{itemize}
\tightlist
\item
 Whitney, H. ``The Work of John Milnor.'' ICM 1962 Laudatio.
\item
 Bass, H. ``John Milnor: The Algebra Behind the Topology.'' \emph{Abel
 Prize Laudatio}, 2011.
\item
 Abate, M. ``Exotic Spheres and John Milnor.'' In \emph{Imagine Math
 2}, Springer, 2013.
\end{itemize}

\begin{center}\rule{0.5\linewidth}{0.5pt}\end{center}

\emph{(End of report)}
